\documentclass[10pt,letterpaper]{article}
\usepackage{cornell-notes}

\title{Clause 27.10.14: Greatest Common Divisor}
\author{}
\date{}
\setDocTitle{Clause 27.10.14: Greatest Common Divisor}
\setDocSubtitle{C++ 2024}
\setDocOwner{C++ 2024 Cornell Notes}
\setDocDate{\today}
\setCornellCollection{C++ 2024 Cornell Notes}
\setCornellUnitType{Clause}
\setCornellUnitNumber{27.10.14}
\setCornellUnitTitle{Greatest Common Divisor}
\hypersetup{pdftitle={Clause 27.10.14: Greatest Common Divisor},pdfsubject={Cornell notes},pdfauthor={}}

\begin{document}
\maketitle
\begin{CornellOverviewBox}{Overview}
\textbf{Classification:} Standard library - Normative\\[2pt]
\textbf{Learning objective:} Understand the rules, terminology, and semantic consequences associated with Greatest common divisor.
\end{CornellOverviewBox}\begin{CornellSummaryBox}{Summary}
{\sffamily\bfseries\color{Accent} Summary}\par\vspace{3pt}Section 27.10.14 focuses on Greatest common divisor. Apply the specified interface together with its mandates, effects, result conditions, exception behavior, and complexity constraints. When applying it, preserve the distinction between mandatory requirements, recommendations, permissions, and behavior deliberately left undefined, unspecified, or implementation-defined.
\end{CornellSummaryBox}

\end{document}
